\documentclass[11pt,a4paper]{article}
\usepackage[utf8]{inputenc}
\usepackage[T1]{fontenc}
\usepackage[brazil]{babel}
\usepackage[a4paper,margin=2.3cm]{geometry}
\usepackage{enumitem}
\usepackage{hyperref}
\usepackage{xcolor}
\setlength{\parindent}{0pt}
\setlength{\parskip}{5pt}
\sloppy
\setlength{\emergencystretch}{3em}
\hypersetup{colorlinks=true,urlcolor=blue,linkcolor=black}
\newcommand{\cabecalho}[2]{\begin{center}{\small \textbf{Sprint 2 -- #1}}\\[1mm]\textit{Inventário de Ativos e Vulnerabilidades}\\Matheus Moreli Durão Silva -- 12621CBS217\\Cibersegurança -- Bacharelado -- Noturno\end{center}\hrule\vspace{4mm}}
\newcommand{\bloco}[1]{\par\vspace{2mm}\textbf{#1}\par}
\newcommand{\codigo}[1]{\texttt{#1}}

\begin{document}\cabecalho{S2\_05 -- Regras e Validações}{}
\bloco{Objetivo} Validar cedo evita registrar dados inconsistentes e evita que uma entrada incorreta encerre o programa. A tela orienta a pessoa; a regra de negócio confirma novamente antes de alterar a base.
\bloco{Regras aplicadas} \codigo{validar\_texto} exige campos obrigatórios; ID deve ser inteiro positivo e único; criticidade, prioridade e status devem pertencer às listas padronizadas; CVE deve seguir \codigo{CVE-AAAA-NÚMERO}; CVSS é convertido para número e limitado entre 0.0 e 10.0; vulnerabilidade somente é anexada a ativo existente.
\bloco{Exemplo} Ao receber \codigo{"11"}, \codigo{validar\_cvss} levanta \codigo{DadosInvalidosError} com mensagem compreensível. A interface captura o erro, mostra a mensagem e mantém o programa aberto. Nenhum \textit{stack trace} nem dado interno é exibido.
\bloco{Como testar} Executar os casos de CVE e CVSS inválidos, ID duplicado e ativo inexistente no arquivo \codigo{test\_sprint2.py}. Cada um espera explicitamente a exceção de regra.
\bloco{Checklist atendido} Obrigatórios verificados, limites testados, execução contínua, erro seguro e regras críticas cobertas.
\end{document}
